\documentclass[11pt]{article}
\usepackage{amsmath}
\title{A timely range can still use an inconsistent pair of frames}
\date{26 November 2026}
\begin{document}
\maketitle
\paragraph{Native boundary and model law.}
The pinned Navigation2 collision-monitor source checks sample age.
With \texttt{base\_shift\_correction} it requests a transform to the current
base time from the scan's acquisition time through its configured fixed
frame. The manager passes those frames and stamps to the TF buffer.
Without base-shift correction the source requests a latest transform
between sensor and base. The selected scan and source methods execute
unchanged against a declared, exact two-sample transform provider, not
ROS TF buffer execution or observed calibration behavior. The modeled
stationary obstacle is observed on an aligned ray.

\paragraph{The joint evidence condition.}
Let $P_i$ take coordinates in the base at time $i$ to a stable reference,
and let $E$ be the fixed sensor-to-base extrinsics. The correct current
coordinates of a stationary point seen at $s$ are $P_t^{-1}P_s E p$.
If the reported poses instead have left-multipliers $G_s,G_t$, the
result is $P_t^{-1}G_t^{-1}G_sP_s E p$. The same unknown multiplier
cancels exactly. Distinct multipliers cannot guarantee equality for
all points, although a particular point may be fixed by their relative
rotation. A common frame label, age threshold, or successful lookup
does not provide this equality. An unchanged generation is only
useful insofar as it warrants the within-generation trajectory law.

\paragraph{A matched contrast.}
The fixture moves the base from 0 to 1 while observing a wall at 4.
The raw delayed range is 4 and true current range is 3. An unchanged
map bias cancels. If it decreases by 2 between stamps, the selected
code returns 5 from the declared transform pair. A declared stable
relative-motion reference gives 3 on the same delayed scan, and a
new scan at the current stamp gives 3. Turning off motion compensation
returns 4, not 3. An ideal rule requiring range above 3.5 would
therefore pass both wrong values; this threshold is part of the
authored comparison, not the native collision policy.

The replay checks 225 transform cases. The alternatives require a
relative reference not affected by that particular change or a fresh
measurement. Neither qualified sensors, timing synchronization, field
behavior, collision avoidance, nor braking dynamics are established.
This is a scoped P06 probe, not an architecture disposition.
\end{document}
